\section*{Capítulo \ref{ch:b1:extent-q-and-k} --- Describir un sistema químico: avance, actividades, Q y K}

\begin{solution}{exo:b1:extent-q-and-k:1}
\omterm{def:b1:extent-q-and-k:variables}{Intensivas}: la temperatura, la \omterm{def:b1:extent-q-and-k:composition}{concentración molar}, la densidad, la
presión y la \omterm{def:b1:extent-q-and-k:composition}{fracción molar}. \omterm{def:b1:extent-q-and-k:variables}{Extensivas}: el volumen, la cantidad de
sustancia y la masa.
\end{solution}

\begin{solution}{exo:b1:extent-q-and-k:2}
$p(\ce{N2}) = 0.7808 \times 1.013 = \qty{0.791}{bar}$, $p(\ce{O2}) =
0.2095 \times 1.013 = \qty{0.212}{bar}$, $p(\ce{Ar}) = 0.0093 \times 1.013
= \qty{0.0094}{bar}$.
\end{solution}

\begin{solution}{exo:b1:extent-q-and-k:3}
$n(\ce{H2}) = 3.0 - 2\xi$, $n(\ce{O2}) = 2.0 - \xi$, $n(\ce{H2O}) =
2\xi$. $\xi_{\max} = \min(3.0/2,\ 2.0/1) = \qty{1.5}{mol}$: el dihidrógeno
es el limitante. Estado final: \ce{H2} 0, \ce{O2} \qty{0.5}{mol}, \ce{H2O}
\qty{3.0}{mol}.
\end{solution}

\begin{solution}{exo:b1:extent-q-and-k:4}
(a) $Q = p_{\ce{CO2}}/p^\circ$. (b) $Q = (c^\circ)^2/([\ce{Ag+}][\ce{Cl-}])$.
(c) $Q = (p_{\ce{NH3}}/p^\circ)^2/[(p_{\ce{N2}}/p^\circ)(p_{\ce{H2}}/p^\circ)^3]$.
(d) $Q = [\ce{NH4+}][\ce{OH-}]/([\ce{NH3}]\,c^\circ)$, pues el agua tiene
\omterm{def:b1:extent-q-and-k:activity}{actividad} 1.
\end{solution}

\begin{solution}{exo:b1:extent-q-and-k:5}
$\mathrm{A \rightleftharpoons C}$ es la suma: $K = 10 \times 0.5 = 5$. $\mathrm{C \rightleftharpoons A}$: $1/5 =
0.2$. $\mathrm{2\,A \rightleftharpoons 2\,C}$: $5^2 = 25$.
\end{solution}

\begin{solution}{exo:b1:extent-q-and-k:6}
$Q = 5.0 \times 10^{-5} < K^\circ$: sentido directo. $Q = 0.10 > K^\circ$:
sentido inverso. Solo reactivos: $Q = 0 < K^\circ$, sentido directo.
\end{solution}

\begin{solution}{exo:b1:extent-q-and-k:7}
$n(\ce{A2}) = 1 - \xi$, $n(\ce{A}) = 2\xi$, total $1 + \xi$. Con $p =
p^\circ$, $Q = \dfrac{(2\xi/(1+\xi))^2}{(1-\xi)/(1+\xi)} = \dfrac{4\xi^2}{1 -
\xi^2}$. $Q = 0.25$ da $4.25\,\xi^2 = 0.25$, $\xi = \qty{0.243}{mol}$.
\end{solution}

\begin{solution}{exo:b1:extent-q-and-k:8}
Con $x = [\ce{C}]$: $x/(0.10 - x)^2 = 100$, de modo que $100x^2 - 21x + 1 = 0$ y
$x = (21 - \sqrt{41})/200 = \qty{0.0730}{mol/L}$ (la otra raíz supera
0.10). $[\ce{A}] = [\ce{B}] = \qty{0.0270}{mol/L}$, $[\ce{C}] =
\qty{0.0730}{mol/L}$, $\tau = 0.73$.
\end{solution}

\begin{solution}{exo:b1:extent-q-and-k:9}
A partir de $n$ de cada uno: $Q = \xi^2/(n - \xi)^2 = K^\circ$, de modo que $\xi/(n - \xi) =
\sqrt{K^\circ}$ y $\tau = \xi/n = \sqrt{K^\circ}/(1 + \sqrt{K^\circ})$.
$\tau = 0.99$ exige $\sqrt{K^\circ} = 99$, $K^\circ \approx \num{9.8e3}$:
del orden de $10^4$, el umbral habitual de una \omterm{def:b1:extent-q-and-k:quantitative}{reacción cuantitativa}.
\end{solution}

\begin{solution}{exo:b1:extent-q-and-k:10}
En el equilibrio, $p_{\ce{CO2}} = 0.20\,p^\circ$, es decir, $n(\ce{CO2}) =
pV/RT = 0.20 \times 10^5 \times 0.0100/(8.314 \times 1100) =
\qty{0.0219}{mol}$. Con \qty{0.010}{mol} de carbonato no puede
alcanzarse: se descompone todo, $p_{\ce{CO2}} = 0.010 \times 8.314 \times
1100/0.0100 = \qty{9.15e3}{Pa} = \qty{0.091}{bar}$, $Q < K^\circ$, y el
estado final (\qty{0.010}{mol} de \ce{CaO}, \qty{0.010}{mol} de \ce{CO2},
nada de \ce{CaCO3}) no es un equilibrio. Con \qty{0.050}{mol}:
equilibrio, \qty{0.0219}{mol} de \ce{CO2} y de \ce{CaO}, y quedan
\qty{0.0281}{mol} de \ce{CaCO3}.
\end{solution}

\begin{solution}{exo:b1:extent-q-and-k:11}
En el equilibrio, $n(\ce{B}) = 2\,n(\ce{A})$ y $n(\ce{C}) = 3\,n(\ce{A})$
(el volumen se simplifica). Entonces, $6\,n(\ce{A}) = 1.0$: $n(\ce{A}) =
\qty{0.167}{mol}$, $n(\ce{B}) = \qty{0.333}{mol}$, $n(\ce{C}) =
\qty{0.500}{mol}$.
\end{solution}

\begin{solution}{exo:b1:extent-q-and-k:12}
Cantidades iguales: $\xi^2/(0.50 - \xi)^2 = 4$, $\xi/(0.50 - \xi) = 2$,
$\xi = \qty{0.333}{mol}$, \omterm{def:b1:extent-q-and-k:fractional-extent}{rendimiento} $0.333/0.50 = 67\,\%$. Con
\qty{1.0}{mol} de alcohol: $\xi^2 = 4(0.50 - \xi)(1.0 - \xi)$,
$3\xi^2 - 6\xi + 2 = 0$,
$\xi = (6 - \sqrt{12})/6 = \qty{0.423}{mol}$, \omterm{def:b1:extent-q-and-k:fractional-extent}{rendimiento} del 85\,\%: un
exceso de un reactivo aumenta el \omterm{def:b1:extent-q-and-k:fractional-extent}{rendimiento} respecto del otro.
\end{solution}

\begin{solution}{pb:b1:extent-q-and-k:1}
\textbf{1.} $x(\ce{CH4}) = 1.00/4.00 = 0.250$; $x(\ce{H2O}) = 0.750$.
\textbf{2.} \qty{7.5}{bar} y \qty{22.5}{bar}.
\textbf{3.} Masas: 16.04 y \qty{54.06}{g}; $w(\ce{CH4}) = 0.229$,
$w(\ce{H2O}) = 0.771$.
\textbf{4.} Las \omterm{def:b1:extent-q-and-k:composition}{fracciones molares}, las \omterm{def:b1:extent-q-and-k:composition}{fracciones másicas} y las \omterm{def:b1:extent-q-and-k:composition}{presiones
parciales} (y la presión total); las cantidades son \omterm{def:b1:extent-q-and-k:variables}{extensivas}.
\textbf{5.} $a = p_i/p^\circ$: 7.5 para el metano y 22.5 para el vapor.
\textbf{6.} \ce{CH4} $1.00 - \xi$, \ce{H2O} $3.00 - \xi$, \ce{CO} $\xi$,
\ce{H2} $3\xi$, total $4.00 + 2\xi$.
\textbf{7.} $\xi_{\max} = \qty{1.00}{mol}$; el metano es el limitante.
\textbf{8.} \ce{CH4} 0, \ce{H2O} \qty{2.00}{mol}, \ce{CO} \qty{1.00}{mol},
\ce{H2} \qty{3.00}{mol}.
\textbf{9.} \qty{6.00}{mol}, frente a \qty{4.00}{mol} en la alimentación:
la reacción forma cuatro moléculas a partir de dos.
\textbf{10.} \ce{H2O} 0.333, \ce{CO} 0.167, \ce{H2} 0.500.
\textbf{11.} Para asegurarse de que reacciona todo el metano y dejar vapor
para la reacción de desplazamiento, cuyo cociente rebaja.
\textbf{12.} \ce{CO} $1 - \xi$, \ce{H2O} $2 - \xi$, \ce{CO2} $\xi$, \ce{H2}
$3 + \xi$, total 6.00 en cualquier avance.
\textbf{13.} $Q = \dfrac{(p_{\ce{CO2}}/p^\circ)(p_{\ce{H2}}/p^\circ)}
{(p_{\ce{CO}}/p^\circ)(p_{\ce{H2O}}/p^\circ)}$ con $p_i = (n_i/6)p$: los
factores $p/(6p^\circ)$ se simplifican entre numerador y denominador (dos
moléculas de gas en cada lado), y queda la expresión dada.
\textbf{14.} $\xi = 0$: $Q = 0 < 4.0$, sentido directo.
\textbf{15.} $\xi(3 + \xi) = 4(1 - \xi)(2 - \xi)$, es decir, $3\xi^2 - 15\xi +
8 = 0$: $\xi = (15 - \sqrt{129})/6 = \qty{0.607}{mol}$; la otra raíz,
4.39, supera $\xi_{\max} = 1$.
\textbf{16.} \ce{CO} \qty{0.393}{mol}, \ce{H2O} \qty{1.393}{mol},
\ce{CO2} \qty{0.607}{mol}, \ce{H2} \qty{3.607}{mol}.
\textbf{17.} $\tau = 0.607/1.00 = 0.607$.
\textbf{18.} $\xi(3 + \xi) = 4(1 - \xi)(5 - \xi)$, $3\xi^2 - 27\xi + 20 =
0$, $\xi = \qty{0.814}{mol}$: más vapor, más conversión.
\textbf{19.} $K^\circ = (0.99 \times 3.99)/(0.01 \times 1.01) = 391$.
\textbf{20.} $Q(\xi)$ es creciente, de modo que una $K^\circ$ mayor se
alcanza con un avance mayor. La constante de esta reacción es mayor a
temperatura más baja, donde la reacción es lenta: un primer reactor
trabaja caliente y rápido, y un segundo, más frío, completa la conversión
(la variación de $K^\circ$ con la temperatura se trata en el volumen del
segundo año).
\textbf{21.} $3.607/(6.00 - 1.393) = 3.607/4.607 = 0.783$.
\textbf{22.} \qty{0.393}{mol} de \ce{CO} y \qty{0.607}{mol} de \ce{CO2}.
\textbf{23.} $x(\ce{H2}) = 3.607/6.00 = \textbf{0.60}$: seis de cada diez
moléculas que salen del reactor de desplazamiento son de hidrógeno.
\end{solution}
